% =====================================================================
%  WHERE THE NEW FIGURES (AND TEXT) GO IN POLISHingEuclid.tex
%
%  Line numbers refer to "POLISHingEuclid (2).tex" (the 11-09 export).
%  If they have moved, search for the FIND text instead.
%
%  Upload first (keep the folder names):
%    images/star_population.pdf          new star figure
%    images/thumbs/lr_rgb.png            diagram thumbnails; your preamble
%    images/thumbs/sr_rgb.png            already points \mdthumbs here
%  and append references_new.bib to your references.bib.
%
%  You can then delete images/fig03_stellar_population_calibration.png.
% =====================================================================


% ---------------------------------------------------------------------
%  1. STAR FIGURE
%     Section: Methods > Training data > Distributions & Transformations
%
%  FIND (lines 336-341), right after the galaxy figure:
%
%    \begin{figure}
%        \centering
%        \includegraphics[width=1\linewidth]{images/fig03_stellar_population_calibration.png}
%        \caption{Star population}
%        \label{fig:staaaars}
%    \end{figure}
%
%  REPLACE WITH the block below.
%  It is figure* now: the plate is one row of four panels, as wide as the
%  page (7.1 in, drawn at 1:1, so the fonts match the text).
%  The label is unchanged, so \ref{fig:staaaars} keeps working.
%  A figure* lands at the top of a page at the earliest one page after
%  where it sits in the source, so keep it close to the star paragraphs.
% ---------------------------------------------------------------------
\begin{figure*}
    \centering
    \includegraphics[width=\textwidth]{images/star_population.pdf}
    \caption{Star population of the training scenes, fitted to Q1 alone.
    (a) VIS counts of Q1 point sources weighted by their photometric-redshift star probability (blue), the straight law fitted in the shaded range (solid) and its extrapolation (dashed).
    \textit{Gaia} DR3 $G$-band counts in the Q1 deep fields (open squares) are shown for comparison only.
    (b--d) VIS$-Y_{\rm E}$, $Y_{\rm E}-J_{\rm E}$ and $J_{\rm E}-H_{\rm E}$ against VIS magnitude for the $4{,}042$ Q1 stars with S/N\,$\geq5$ in all bands (blue) and for model stars drawn at the same magnitudes and re-observed with each star's Q1 errors (orange).}
    \label{fig:staaaars}
\end{figure*}


% ---------------------------------------------------------------------
%  2. STAR TEXT (for completeness)
%
%  INSERT the two star paragraphs from model_and_stars_options.tex
%  (lines 11-71: "STARS, paragraph 1" and "STARS, paragraph 2")
%  AFTER line 325:
%
%    We use extreme deconvolution method \citep{Bovy2011} to remove the
%    measurement noise from each region.
%
%  i.e. just above your "MORE MORE MORE PLOTS!" note.
% ---------------------------------------------------------------------


% ---------------------------------------------------------------------
%  3. MODEL SUBSECTION
%
%  FIND (lines 349-357):
%
%    \subsection{Training}
%    \subsubsection{Choices}
%    \begin{itemize}
%        \item asinh knee
%        \item loss
%        \item training steps
%        \item depth
%    \end{itemize}
%    \subsubsection{Combiner model}
%
%  REPLACE WITH \subsection{Model}\label{sec:model} and its four paragraphs
%  from model_and_stars_options.tex (lines 86-204).
%  Then move the architecture figure* (section 4 below) so it sits right
%  after the first Model paragraph (architecture), not after the last one.
% ---------------------------------------------------------------------


% ---------------------------------------------------------------------
%  4. ARCHITECTURE FIGURE (the TikZ figure*, lines 359-445)
%
%  Option A (recommended: three edits, your TikZ otherwise untouched)
%
%  (i)  line 390, FIND:
%         \node[mdsum] (sum) at (12.55,0) {$+$};
%       REPLACE WITH:
\node[mdsum] (sum) at (12.75,0) {$+$};
%
%  (ii) line 392, FIND:
%         \node[mdbox,draw=mdKnee!80!black,fill=mdKnee!18,minimum height=\hL cm] (knee) at (13.95,0)
%       REPLACE WITH:
\node[mdbox,draw=mdKnee!80!black,fill=mdKnee!18,minimum height=\hL cm] (knee) at (14.45,0)
%       (these two moves open up the arrow from the sum into the knee box,
%       which was squashed to nothing)
%
%  (iii) line 443, FIND:
%         \caption{...}
%       REPLACE WITH:
\caption{The reconstruction network.
(a) The four-band input on the $0.1''$ grid is stretched with $\mathrm{arcsinh}$ at eight knees, giving $32$ channels.
The shared trunk (a $3\times3$ convolution, $16$ residual blocks, a second $3\times3$ convolution and a pixel shuffle) predicts a correction that is added to a per-band path, in which each band is upsampled from its own eight channels by a $5\times5$ convolution and a pixel shuffle.
A learned knee $k_b$ for each band converts the sum to electrons on the $0.05''$ grid.
Numbers above the blocks are channel counts; gray blocks lie on the input grid and green blocks on the output grid.
(b) One residual block: a $1\times1$ convolution widens $32$ channels to $192$ before the ReLU, a linear $1\times1$ convolution narrows them to $25$, and a $3\times3$ convolution returns $32$ channels, which are added to the block input.
The thumbnails show a Q1 lens candidate and its reconstruction (R\,=\,$H_{\rm E}$, G\,=\,mean of $Y_{\rm E}$ and $J_{\rm E}$, B\,=\,VIS).}
%
%  Option B: replace the whole block (lines 359-445, \begin{figure*} ...
%  \end{figure*}) with architecture_figure.tex. Same result; it only
%  differs from yours in comments and the two coordinates above.
%
%  Option C: use the prebuilt PDF instead of TikZ. Replace lines 359-445
%  with the block below (keep the caption from (iii)). The PDF is drawn at
%  the page width (512.8 pt), so it is shown at 1:1. You can then delete
%  the "TIKZ PREAMBLE" (from %%%% TIKZ PREAMBLE down to the end of the
%  \mdimstack definition) and the thumbs. Downside: the diagram no longer
%  changes when you edit the TikZ.
%
%    \begin{figure*}
%        \centering
%        \includegraphics[width=\textwidth]{images/model_diagram.pdf}
%        \caption{...}   % caption from (iii)
%        \label{fig:architecture}
%    \end{figure*}
% ---------------------------------------------------------------------
